\documentclass[10pt,a4paper]{amsart}
\usepackage[margin=19mm]{geometry}
\usepackage{amsmath,amssymb,mathtools}
\usepackage[scr=rsfs,cal=euler]{mathalfa}
\usepackage{xfrac}
\usepackage[unicode,hidelinks,bookmarks=false]{hyperref}
\newcommand{\ie}{i.e.\ }
\newcommand{\cf}{cf.\ }
\newcommand{\resp}{respectively\ }
\newcommand{\Subsection}{\subsection}
\providecommand{\nobreakdash}{\nobreak\hbox{-}}
\setlength{\emergencystretch}{2em}
\multlinegap0pt
\usepackage{amssymb}
\usepackage{mathtools}
\usepackage{float}
\usepackage{tikz}
\usepackage{xcolor}
\newtheorem{lemma}{Lemma}
\title{Induced-saturated graphs exist for even cycles}
\author{Ilkyoo Choi}
\date{Source: arXiv:2608.24202v1 (CC BY 4.0)}
\begin{document}
\maketitle

\noindent\emph{Source and licence.} Induced-saturated graphs exist for even cycles. Version arXiv:2608.24202v1 (CC BY 4.0), distributed by arXiv under the Creative Commons Attribution 4.0 International licence, http://creativecommons.org/licenses/by/4.0/, as recorded in the arXiv licence metadata for this exact version and shown on the abstract page https://arxiv.org/abs/2608.24202v1. Full licence notice: \texttt{licenses/arxiv-2608.24202-CC-BY-4.0.txt}.\par\smallskip

\noindent\textit{Editorial scope.} Counted unit: the article's lemma lem:routing (source0.tex lines 987--989). The article's complete original proof of it is delivered with it (source0.tex lines 990--1023, one \texttt{proof} environment of 34 lines). No mathematics is written, repaired or re-derived: the statement and the proof are the article's own lines, in the article's own notation, and no retained line is substituted or re-flowed.  No further source-local context is needed: every cross-reference of every retained line resolves inside the two delivered spans.  Nothing was omitted from any retained span, so the package carries no omission record.  No retained line contains a citation, so the wrapper carries no bibliography and no bibliography entry.  No retained line contains a figure environment, an image-inclusion command or drawing code, so no image is delivered and the package declares no asset dependency.  Editorial formatting only, in addition: an AMS \texttt{amsart} preamble that restates the article's own declarations and macros used by the retained lines, namely the environments \texttt{lemma} on separate counters, together with the standard packages those lines need.  The retained environments print in this document as Lemma 1 in order of appearance, whereas the article prints the same items with the numbering of its own full document.  The complete native source of the article as distributed in the arXiv e-print for this exact version is bundled unchanged as \texttt{sources/208456/source0.tex}.
\begin{lemma}\label{lem:routing}
Whenever defined, each sequence $\sigma_{0^{10}}, \sigma_{0^{11}}, \sigma_{1^{10}}, \sigma^{\mathrm{even}}_{d}, \sigma^{\mathrm{odd}}_{d}, \sigma_{3}, \sigma_{4}, \sigma_{5}$ is an induced path on $2q+2$ vertices. 
\end{lemma}

\begin{proof}
It is easy to check that each sequence has $2q+2$ vertices.
Since edges of $G_q$ are only within blocks and in connections, we only need to check edges and non-edges within a block and in connections. 

Consider $\sigma_{0^{10}}$.
It has only one matching edge, which is $0^{10}0^{11}$.
It also has exactly one cross edge in every connection. 

Consider $\sigma_{0^{11}}$.
It has only one matching edge, which is $q^{00}q^{01}$.  
It also has exactly two cross edges in the $i$th connection where $i\in\{q, \ldots, 2q-1\}$. 

Consider $\sigma_{1^{10}}$.
It has no matching edges. 
It has exactly one cross edge in every connection, except for the $0$th connection where it has exactly two cross edges. 

Consider  $\sigma_d^{\mathrm{even}}$.
It has exactly one matching edge $r^{00}r^{01}$. 
It also has exactly one cross edge in the $i$th connection where $i\in\{0, \ldots, d-1\}$ and exactly two cross edges in the $i$th connection where $i\in\{r, \ldots, 2q-1\}$. 

Consider  $\sigma_d^{\mathrm{odd}}$.
It has exactly two matching edges $r^{00}r^{01}$ and $2^{00}2^{01}$. 
It also has exactly one cross edge in the $i$th connection where $i\in\{0, \ldots, d-1\}\setminus\{2\}$, exactly two cross edges in the $i$th connection where $i\in\{r, \ldots, 2q-1\}$, and exactly three cross edges in the 2nd connection. 

Each of $\sigma_3, \sigma_4, \sigma_5$ has two matching edges. 
$\sigma_3$ has exactly one cross edge in the $i$th connection where $i\in\{0,1, 2\}$ and exactly two cross edges in the 3rd connection. 
$\sigma_4$ has exactly one cross edge in the $1$st and 2nd connection, exactly two cross edges in the 3rd connection, and exactly three cross edges in the 0th connection. 
$\sigma_5$ has exactly one cross edge in the $0$th, 1st, 3rd, 4th connection, exactly two cross edges in the 5th connection, and exactly three cross edges in the 2nd connection. 

For every sequence $\sigma$, consecutive vertices in $\sigma$ are adjacent.
Moreover, since edges of $G_q$ are either  within a block or between consecutive blocks, the matching edges and cross edges listed above are all the edges induced by $\sigma$.
In each case, there are exactly  $2q+1$ edges.
Hence each sequence is an induced path on $2q+2$ vertices.
\end{proof}

\end{document}
